%!TEX root = document.tex

\section{安装和配置}

\subsection{系统要求}

md-highlighter 是一个VSCode扩展。本项目没有单独的后端程序，也不需要数据库。实际使用时，主要要求是VSCode正常加载扩展配置和TextMate语法文件。

\begin{table}[h]
\centering
\caption{md-highlighter 系统要求}
\begin{tabular}{p{3cm}p{10cm}}
\toprule
\textbf{要求项目} & \textbf{规格} \\
\midrule
操作系统 & Windows、macOS、Linux，和VSCode支持范围一致 \\
VSCode版本 & 1.70.0及以上版本，对应\lstinline|package.json|中的\lstinline|engines.vscode|配置 \\
硬件要求 & 无单独硬件要求，通常与VSCode编辑普通文本文件时的要求一致 \\
运行依赖 & 无额外服务端依赖，主要由VSCode加载本扩展中的JSON配置文件 \\
网络要求 & 安装扩展时可能需要网络；本地使用语法高亮时不依赖网络连接 \\
\bottomrule
\end{tabular}
\end{table}

\subsection{软件特性}

当前版本在\lstinline|package.json|中注册了23个语言ID，对应23个语法定义文件，范围包括NAMD/CHARMM、Amber/AmberTools、Gromacs、小分子结构文件以及PLUMED输入文件。

\subsubsection{颜色编码方式}

该扩展本身并不直接定义每一种颜色，而是为文本赋予TextMate作用域，最终由当前主题决定显示效果。常见的作用域用法如下：

\begin{itemize}
    \item \textbf{名称类字段}：原子名、残基名、分段名称等，一般使用\lstinline|entity.name.*|或相近作用域
    \item \textbf{数值类字段}：坐标、电荷、质量、一般参数值，常见为 \lstinline|constant.numeric.*|
    \item \textbf{类型类字段}：原子类型、参数类型、格式类型等，常见为 \lstinline|support.type.*|
    \item \textbf{注释和记录头}：注释行、说明行、记录关键字，会根据具体文件使用 \lstinline|comment.*|、\lstinline|keyword.*| 等作用域
\end{itemize}

这种实现方式下，语法文件主要负责“分段”和“标记”，无需在扩展中另外维护颜色表。

\subsubsection{主题兼容性}

在\lstinline|Atom One Light|主题下已经测试显示效果。对于其他主题，本扩展仍按VSCode的作用域取色机制显示；但不同主题之间的颜色差异较大，实际显示效果以用户本地主题为准。

如需手动调整某一类字段的颜色，可在\lstinline|settings.json|中覆盖指定作用域。例如：

\begin{lstlisting}[style=blockstyle,language=json]
{
    "editor.tokenColorCustomizations": {
        "textMateRules": [
            {
                "scope": "entity.name.tag.atom-name",
                "settings": {
                    "foreground": "#0066CC",
                    "fontStyle": "bold"
                }
            }
        ]
    }
}
\end{lstlisting}

\subsection{安装方法}

\subsubsection{从Visual Studio Marketplace安装}

若该版本已经发布到扩展市场，可直接在VSCode中搜索安装。操作步骤如下：
\begin{itemize}
    \item 打开VSCode
    \item 进入扩展视图，或者按 \texttt{Ctrl+Shift+X}
    \item 搜索 \lstinline|md-highlighter|
    \item 选择发布者为 \lstinline|gxf1212| 的扩展并安装
\end{itemize}

\subsubsection{从VSIX文件安装}

若从仓库发布页获取VSIX包，也可手动安装：
\begin{itemize}
    \item 打开VSCode
    \item 按 \texttt{Ctrl+Shift+P} 打开命令面板
    \item 输入 \lstinline|Extensions: Install from VSIX...|
    \item 选择对应的VSIX文件
    \item 等待VSCode完成安装
\end{itemize}

\subsubsection{从命令行安装}

若本机已经能使用\lstinline|code|命令，也可直接执行：
\begin{lstlisting}[style=blockstyle,language=bash]
code --install-extension gxf1212.md-highlighter
\end{lstlisting}

\subsection{验证安装}

安装完成后，可打开一个已经注册的文件类型，查看语言模式和高亮是否生效。可按以下方法确认：

\begin{itemize}
    \item 打开一个 \lstinline|.pdb|、\lstinline|.rtf|、\lstinline|.gro| 或 \lstinline|.plumed.dat| 文件
    \item 查看VSCode右下角的语言模式是否切换到本扩展注册的语言ID
    \item 查看文件内容是否已按字段出现分段高亮，而不是整段普通文本
    \item 在扩展列表中确认 \lstinline|md-highlighter| 处于已启用状态
\end{itemize}

\subsection{配置选项}

当前版本没有自定义设置项，安装后即可使用。与显示有关的调整主要通过VSCode自身的主题和文件关联能力完成，而不是由本扩展单独提供配置面板。

\subsubsection{主题适配}

主题适配主要取决于当前主题对TextMate作用域的解释方式。\lstinline|Atom One Light|主题下已验证显示效果。若用户本地使用其他主题，仍可显示高亮，但具体颜色深浅和对比度取决于对应主题本身的配色方案。

\subsubsection{文件关联}

扩展在\lstinline|package.json|里注册了扩展名和部分文件名模式。大多数情况下，打开文件时会自动进入对应语言模式。比如：
\begin{itemize}
    \item \lstinline|.prmtop| 和 \lstinline|.parm7| 会进入 \lstinline|prmtop| 模式
    \item \lstinline|.prepin| 和 \lstinline|.prepi| 会进入 \lstinline|prepin| 模式
    \item \lstinline|leaprc.test| 会按 \lstinline|in| 语言模式处理
    \item \lstinline|.plumed.dat| 会进入 \lstinline|plumed| 模式
\end{itemize}

如果文件后缀不标准，或者VSCode未自动识别，也可手动切换语言模式，再选择本扩展注册的条目。

\subsection{更新和维护}

\subsubsection{更新方式}

如果扩展通过市场安装，更新方式和普通VSCode扩展一致，通常由编辑器自动处理。如果使用的是本地VSIX包，则需要在版本更新后重新安装新版包。

\subsubsection{版本历史}

\lstinline|CHANGELOG.md|记录了0.0.1到0.0.7的变更。当前版本0.0.7的明确增量是增加了\lstinline|vsd|文件支持。更早几个版本则陆续补充了\lstinline|parm7|、\lstinline|rst7|、\lstinline|plumed|、\lstinline|prepin|、\lstinline|frcmod|、\lstinline|atp|、\lstinline|rtp|、\lstinline|hdb|、\lstinline|r2b|、\lstinline|tdb|、\lstinline|arn|等格式。

\subsubsection{问题反馈}

保留了两个反馈渠道：
\begin{itemize}
    \item GitHub Issues：\url{https://github.com/gxf1212/md-highlighter/issues}
    \item 邮件：\href{mailto:gxf1212@zju.edu.cn}{gxf1212@zju.edu.cn}
\end{itemize}

\subsection{卸载}

卸载方式和普通VSCode扩展一样：
\begin{itemize}
    \item 打开扩展管理界面
    \item 找到 \lstinline|md-highlighter|
    \item 点击卸载
    \item 如有需要，重启VSCode
\end{itemize}

卸载后只是移除语言注册和语法高亮配置，不会改动用户原来的文本文件。
